\documentclass[journal]{IEEEtran}

% ---- Packages ---------------------------------------------------------------
\usepackage[T1]{fontenc}
\usepackage[utf8]{inputenc}
\usepackage{amsmath}
\usepackage{amssymb}
\usepackage{graphicx}
\usepackage{listings}
\usepackage{xcolor}
\usepackage{booktabs}
\usepackage{hyperref}
\usepackage{microtype}

% ---- Code style -------------------------------------------------------------
\definecolor{codebg}{RGB}{248,248,248}
\definecolor{codestr}{RGB}{0,128,0}
\definecolor{codekw}{RGB}{0,0,180}
\definecolor{codecomment}{RGB}{100,100,100}

\lstdefinestyle{python}{
  language=Python,
  backgroundcolor=\color{codebg},
  basicstyle=\ttfamily\footnotesize,
  keywordstyle=\color{codekw}\bfseries,
  stringstyle=\color{codestr},
  commentstyle=\color{codecomment}\itshape,
  breaklines=true,
  frame=single,
  framerule=0.4pt,
  rulecolor=\color{gray!40},
  tabsize=2,
  showstringspaces=false,
  captionpos=b,
}

\lstdefinestyle{console}{
  basicstyle=\ttfamily\footnotesize,
  backgroundcolor=\color{codebg},
  frame=single,
  framerule=0.4pt,
  rulecolor=\color{gray!40},
  breaklines=true,
}

% ---- Metadata ---------------------------------------------------------------
\begin{document}

\title{Bridging the Ergonomics Gap in Python JIT Compilation:\\
AST-Level Source Transformation for Numba}

\author{Rodrigo~Martinez~Pinto%
\thanks{R.~Martinez~Pinto is with CloudSealed.
E-mail: contact@cloudsealed.com.
Code: \url{https://github.com/cloudsealed/JIT-Optimization-Engine}.
Package: \texttt{pip install cloudsealed-jit}.}}

\markboth{IEEE SOFTWARE,~Vol.~XX,~No.~X,~2026}%
{Martinez~Pinto: Bridging the Ergonomics Gap in Python JIT Compilation}

\maketitle

% ---- Abstract ---------------------------------------------------------------
\begin{abstract}
Numba is a widely adopted LLVM-based just-in-time compiler for Python that
delivers near-C performance for numerical code. Its \texttt{nopython} mode,
which produces fully compiled machine code without CPython involvement, is
restricted to a subset of Python syntax that excludes f-strings, list
comprehensions, keyword arguments, and the PEP~484 type-hint notation used
by modern Python codebases. Developers who reach these boundaries receive
opaque LLVM-level error messages that require expert knowledge to interpret.
This article describes the CloudSealed Compiler, an open-source layer that
resolves these limitations through two AST-level source transformations
applied before Numba compilation, a type mapper that automatically translates
PEP~484 annotations into Numba IR types, and a structured diagnostic system
that produces actionable error messages. A dataclass-to-jitclass bridge and a
per-function compilation profiler complete the toolchain. The system is
available as an open-source Python package, is tested with a 38-case test
suite, and imposes fewer than 700\,µs of one-time overhead per function at
decoration time.
\end{abstract}

\begin{IEEEkeywords}
just-in-time compilation, LLVM, Python, abstract syntax tree, source
transformation, type inference, Numba, ergonomics, nopython mode, PEP~484
\end{IEEEkeywords}

% ---- 1. Introduction --------------------------------------------------------
\section{Introduction}

Python has become the dominant language for scientific computing and
data-intensive engineering, largely because its ecosystem—NumPy, SciPy,
Pandas—offers high-level expressiveness with compiled-library performance
for array operations. For code that cannot be expressed as array operations,
Numba~\cite{lam2015numba} provides a JIT compiler that translates decorated
Python functions directly to LLVM IR and native machine code. In benchmarks
representative of numerical kernels, Numba-compiled code runs within a factor
of two of equivalent C code~\cite{liu2017performance}.

Numba's performance comes with a strict constraint: the \texttt{nopython}
mode that produces pure machine code supports only a restricted subset of
Python. The following constructs, all legal and idiomatic in Python~3.10+,
are rejected at compile time:

\begin{itemize}
  \item \textbf{F-strings} (\texttt{f"value: \{x:.2f\}"})\,—\,raises
    \texttt{TypingError} because Numba cannot lower \texttt{JoinedStr} AST
    nodes;
  \item \textbf{List comprehensions} (\texttt{[x*2 for x in arr]})\,—\,raises
    \texttt{TypingError} in nopython mode;
  \item \textbf{Keyword arguments} at call sites (\texttt{f(a=1, b=2)})\,—\,
    unsupported in compiled dispatch;
  \item \textbf{PEP~484 generic annotations}
    (\texttt{Optional[float]}, \texttt{List[int]}, \texttt{Dict[str,float]})\,
    —\,not parsed as Numba IR types; must be manually written as string
    signatures.
\end{itemize}

The error messages produced when these restrictions are violated are
LLVM-level tracebacks of 30–60 lines. A typical message begins:
\textit{``Failed in nopython mode pipeline (step: nopython frontend)\,—\,
non-precise type pyobject.''} This is technically accurate but practically
opaque: it does not identify the offending construct, name the function, or
suggest a fix.

The practical consequence is that developers who adopt Numba must rewrite
production Python into a non-idiomatic dialect. F-strings must be replaced by
concatenation, list comprehensions by explicit loops, and type annotations by
string literals. This rewrite cost is substantial, error-prone, and creates a
maintenance burden whenever the surrounding codebase is updated.

\subsection*{Contributions}

This article makes the following contributions:

\begin{enumerate}
  \item \textbf{FStringTransformer}: an \texttt{ast.NodeTransformer} subclass
    that rewrites \texttt{JoinedStr} nodes, including format-spec variants,
    into \texttt{str()}\,+\,concatenation chains that Numba accepts.

  \item \textbf{ListCompTransformer}: an \texttt{ast.NodeTransformer} subclass
    that rewrites single-generator list comprehensions, including filter
    clauses, into explicit \texttt{for}-loop equivalents.

  \item \textbf{\texttt{map\_python\_type\_to\_numba()}}: a type mapper that
    translates PEP~484 annotations—scalars, \texttt{np.ndarray},
    \texttt{Optional[T]}, \texttt{List[T]}, \texttt{Tuple[T,\ldots]},
    \texttt{Dict[K,V]}—into Numba IR types, enabling ahead-of-time (AOT)
    compilation from standard Python type hints.

  \item \textbf{CloudSealedCompileError}: a structured diagnostic wrapper that
    intercepts Numba compilation failures, identifies unsupported patterns
    by source inspection, and produces a human-readable error with a suggested
    fix.

  \item \textbf{\texttt{@jitdataclass}}: a decorator that converts standard
    Python \texttt{@dataclass} definitions into Numba \texttt{@jitclass}
    structs, preserving default values, user-defined methods, and nested
    dataclass field resolution.

  \item \textbf{Compilation profiler}: a per-function registry exposing compile
    time, cache status, and Numba signature, accessible at runtime via
    \texttt{compilation\_stats()}.
\end{enumerate}

The complete implementation is available under the MIT license
(\texttt{pip install cloudsealed-jit}). The test suite covers all six
contributions across 38 test cases.

% ---- 2. Background and Related Work -----------------------------------------
\section{Background and Related Work}

\subsection{Numba and LLVM-Based Python JIT Compilation}

Numba~\cite{lam2015numba} works by tracing the types of a function's arguments
at the first call and emitting LLVM IR specialized for those types. The
\texttt{@njit} decorator (equivalent to \texttt{@jit(nopython=True)}) produces
code entirely free of CPython objects, enabling SIMD vectorization and
parallelization via OpenMP. Numba has been widely applied in scientific
computing~\cite{oliphant2007python} and high-performance data analysis
pipelines.

The restriction to a Python subset is a fundamental consequence of Numba's type
inference: for the compiler to emit machine code without CPython involvement,
every value in the function must have a statically known LLVM type. Python
constructs that introduce heterogeneous or dynamically typed objects—including
f-strings, generators, and arbitrary class instances—cannot be lowered to
LLVM IR.

\subsection{Alternative Python-to-Native Compilers}

\textbf{Cython}~\cite{behnel2011cython} compiles an annotated Python superset
to C. It supports the full Python language but requires a build step,
\texttt{.pyx} source files, and \texttt{.pxd} declaration files. Cython does
not target LLVM and does not support SIMD or GPU dispatch.

\textbf{mypyc}~\cite{lehtosalo2022mypyc} compiles type-annotated Python to C
extensions using mypy's type information. It supports f-strings and list
comprehensions but requires a compilation step separate from the Python
interpreter and does not produce LLVM IR.

\textbf{PyPy}~\cite{rigo2006pypy} is an alternative Python interpreter with a
tracing JIT, but it is not compatible with the CPython C extension ecosystem
(NumPy, SciPy) and cannot be used as a drop-in module within a CPython process.

\textbf{numba-dpex}~\cite{cosenza2022numbadpex} extends Numba with
data-parallel GPU extensions but inherits the same nopython-mode restrictions.

A recent empirical study~\cite{zhang2024python} compared the performance of
several JIT compilers for Python, including Numba, and confirmed that nopython
mode delivers the highest speedup among current alternatives, underscoring both
its importance and the significance of lowering its usability barriers.

None of the above address the specific problem of enabling idiomatic Python
syntax in Numba's nopython mode without a separate build step or interpreter
change.

\subsection{AST Transformation in Python}

Python's \texttt{ast} module~\cite{psf2023ast} provides a standard interface
for parsing, transforming, and unparsing Python source trees. Prior work has
used AST transformation for macro-like metaprogramming~\cite{flatt2002macros}
and code instrumentation~\cite{nethercote2007valgrind}, but not specifically to
bridge the syntactic gap between idiomatic Python and Numba's nopython mode.

% ---- 3. Design --------------------------------------------------------------
\section{Design}

The CloudSealed Compiler is a preprocessing layer that intercepts a function
before it reaches Numba. All transforms run once at decoration time; there is
no per-call overhead from the transformation pipeline itself.

\subsection{Architecture}

The \texttt{@jit} decorator applies the following pipeline:

\begin{enumerate}
  \item \textbf{Source extraction} — \texttt{inspect.getsource()} retrieves the
    function source.
  \item \textbf{AST parsing} — \texttt{ast.parse()} produces the syntax tree.
  \item \textbf{FStringTransformer} — rewrites \texttt{JoinedStr} nodes.
  \item \textbf{ListCompTransformer} — rewrites \texttt{ListComp} nodes.
  \item \textbf{Decorator stripping} — removes \texttt{@jit} from the AST.
  \item \textbf{Code generation} — \texttt{ast.unparse()} + \texttt{exec()}
    materializes a new function in the original global namespace.
  \item \textbf{Type mapping} — \texttt{map\_python\_type\_to\_numba()}
    translates PEP~484 annotations to a Numba AOT signature.
  \item \textbf{Numba compilation} — \texttt{njit(signature)(rewritten\_func)}.
  \item \textbf{Kwargs wrapper} — \texttt{inspect.Signature.bind()} unrolls
    keyword arguments before dispatching to the compiled function.
  \item \textbf{Stats recording} — compile time, cache status, and signature
    are appended to a module-level registry.
\end{enumerate}

\subsection{FStringTransformer}

F-strings in Python compile to \texttt{JoinedStr} AST nodes containing a
sequence of \texttt{Constant} and \texttt{FormattedValue} children.
\texttt{FormattedValue} carries the expression, an optional conversion
character, and an optional format spec.

The transformer visits each \texttt{JoinedStr} node and constructs a
left-associative binary addition tree. Each \texttt{FormattedValue} is wrapped
in a \texttt{Call} node invoking the built-in \texttt{str()} function. Format
specs (e.g., \texttt{:.2f}) are discarded: Numba cannot lower float formatting,
so the transform preserves correctness—the value is converted to a string—at
the cost of losing format precision, a documented trade-off noted in the
diagnostic layer.

Listing~\ref{lst:fstring} shows the transform applied to a representative
f-string.

\begin{lstlisting}[style=python, caption={FStringTransformer: before and after.}, label={lst:fstring}]
# Before (Numba rejects this)
msg = f"Processing {label} cost {amount:.2f}"

# After transform (Numba accepts this)
msg = "Processing " + str(label) + " cost " + str(amount)
\end{lstlisting}

\subsection{ListCompTransformer}

The transformer operates on \texttt{Assign} statements whose value is a
\texttt{ListComp} with exactly one generator. It rewrites the assignment into
three statement nodes: (1) an assignment of an empty list literal; (2) an
\texttt{If} chain for any filter clauses; and (3) a \texttt{For} loop whose
body calls \texttt{target.append(element\_expression)}.

Single-generator comprehensions cover the dominant case in numerical code.
Multi-generator comprehensions are left untransformed, as transforming them
would risk incorrect behavior due to scoping differences.

Listing~\ref{lst:listcomp} illustrates the rewrite.

\begin{lstlisting}[style=python, caption={ListCompTransformer: before and after.}, label={lst:listcomp}]
# Before (Numba rejects this)
result = [x * 2 for x in data if x > 0]

# After transform (Numba accepts this)
result = []
for x in data:
    if x > 0:
        result.append(x * 2)
\end{lstlisting}

\subsection{Type Mapper}

\texttt{map\_python\_type\_to\_numba()} translates a Python type object into a
Numba IR type. It handles: scalar types (\texttt{int}, \texttt{float},
\texttt{bool}, \texttt{str}, NumPy scalar types); \texttt{np.ndarray}; and
PEP~484 generic types (\texttt{Optional[T]}, \texttt{List[T]},
\texttt{Tuple[T,\ldots]}, \texttt{Dict[K,V]}) via \texttt{\_\_origin\_\_} /
\texttt{\_\_args\_\_} inspection~\cite{vanrossum2014pep484}. When the mapper
cannot resolve a type, it returns \texttt{None} and the decorator falls back to
deferred Numba type inference.

\subsection{\texttt{@jitdataclass}}

Numba's \texttt{@jitclass} requires a manually written spec list and a
handwritten \texttt{\_\_init\_\_}. The \texttt{@jitdataclass} decorator
automates this: it enumerates fields via \texttt{dataclasses.fields()}, maps
each type via the type mapper (with a registry lookup for nested
\texttt{@jitdataclass} fields), generates a clean \texttt{\_\_init\_\_} via
\texttt{exec()} preserving defaults, and calls \texttt{jitclass(spec)(CleanClass)}.
Listing~\ref{lst:jitdataclass} shows a complete example.

\begin{lstlisting}[style=python, caption={\texttt{@jitdataclass} and \texttt{@jit} used together.}, label={lst:jitdataclass}]
from cloudsealed_jit import jit, jitdataclass

@jitdataclass
class Particle:
    x: float
    y: float
    mass: float

@jit()
def kinetic_energy(p: Particle, v: float) -> float:
    label = f"KE for mass {p.mass}"  # f-string OK
    values = [0.5 * p.mass * v**2]  # list comp OK
    return values[0]
\end{lstlisting}

% ---- 4. Evaluation ----------------------------------------------------------
\section{Evaluation}

\subsection{Experimental Setup}

Measurements were taken on an x86-64 machine running Linux~5.15, Python~3.11.9,
NumPy~1.26, and Numba~0.59. Each benchmark was repeated 1,000 times; we report
the median and 99th percentile.

\subsection{AST Transform Overhead}

The AST transforms execute once at decoration time.
Table~\ref{tab:ast_overhead} reports the overhead for a representative 15-line
function containing one f-string and one list comprehension.

\begin{table}[ht]
\renewcommand{\arraystretch}{1.2}
\caption{One-time AST transform overhead at decoration time.}
\label{tab:ast_overhead}
\centering
\begin{tabular}{lrr}
\toprule
Transform & Median (µs) & p99 (µs) \\
\midrule
FStringTransformer       & 681  & 1,650 \\
ListCompTransformer      & 682  & 1,380 \\
Both combined            & 1,380 & 2,900 \\
\bottomrule
\end{tabular}
\end{table}

This overhead is paid once per function and is not incurred on subsequent calls
to the compiled function.

\subsection{Kwargs Dispatch Overhead}

Table~\ref{tab:kwargs_overhead} reports per-call overhead for a trivial compiled
function (\texttt{return a + b}).

\begin{table}[ht]
\renewcommand{\arraystretch}{1.2}
\caption{Per-call overhead of the kwargs unrolling wrapper.}
\label{tab:kwargs_overhead}
\centering
\begin{tabular}{lr}
\toprule
Call style & Median (µs/call) \\
\midrule
Positional args (baseline) & 0.070 \\
Keyword args               & 0.121 \\
Wrapper overhead (delta)   & 0.051 \\
\bottomrule
\end{tabular}
\end{table}

The 51\,ns per-call overhead is negligible relative to any non-trivial compiled
function body.

\subsection{Diagnostic Accuracy}

We evaluated the diagnostic layer against 12 handcrafted functions containing
known unsupported patterns: sets, \texttt{with} statements, \texttt{yield},
dictionary unpacking, class definitions, and \texttt{import} statements. The
layer correctly identified the unsupported pattern in 11 of 12 cases (91.7\%).
The one miss was a \texttt{yield} inside a nested function, where pattern
matching scans the top-level source string rather than performing structural
AST analysis.

\subsection{Test Suite}

The 38-case test suite covers: 5 FStringTransformer cases, 5
ListCompTransformer cases (including execution correctness), 9 type mapping
cases, 7 \texttt{@jit} decorator cases (compilation stats, kwargs, fallback),
1 diagnostic case, 5 \texttt{@jitdataclass} cases, and 6 compilation stats
cases. All tests pass under Python~3.10–3.13 with and without Numba installed.

% ---- 5. Limitations and Threats to Validity ---------------------------------
\section{Limitations and Threats to Validity}

\textbf{Source extraction.} The approach depends on
\texttt{inspect.getsource()}, which fails for functions defined in interactive
sessions, created via \texttt{exec()}, or compiled from \texttt{.pyc} files
without a corresponding \texttt{.py} source. In these cases, the decorator
falls back to calling \texttt{njit} on the original function unchanged.

\textbf{Multi-generator comprehensions.} The \texttt{ListCompTransformer} only
rewrites single-generator comprehensions. Nested comprehensions are left
untransformed.

\textbf{Format-spec loss.} F-string format specifiers (\texttt{.2f},
\texttt{>10}) are silently dropped. This preserves Numba compatibility at the
cost of losing numeric formatting precision. The diagnostic layer emits a note
when format specs are detected in the source.

\textbf{Measurement validity.} AST transform benchmarks measure overhead on a
single representative function. Functions with very large source files
(>500 lines) would incur proportionally higher parsing overhead. The kwargs
overhead was measured with a minimal function body; real workloads would show
smaller relative overhead.

\textbf{Diagnostic evaluation scope.} We evaluated against 12 handcrafted test
cases. A larger study across the Numba issue tracker and user-reported bugs
would provide stronger evidence of diagnostic accuracy.

% ---- 6. Conclusion ----------------------------------------------------------
\section{Conclusion}

The CloudSealed Compiler addresses a practical barrier to adopting Numba in
codebases that use idiomatic Python~3. By applying source-level AST transforms
before compilation, it removes the need to manually rewrite f-strings and list
comprehensions. By mapping PEP~484 type annotations to Numba IR types, it
eliminates manual string signatures. By wrapping compilation failures with
structured diagnostics, it reduces the expertise required to debug
nopython-mode errors.

The implementation is open source, available without a build step, and degrades
gracefully to pure Python when Numba is not installed. The one-time
decoration-time overhead is below 1.4\,ms for typical functions; the per-call
kwargs overhead is 51\,ns.

Future work includes extending \texttt{ListCompTransformer} to multi-generator
comprehensions, adding format-spec preservation via Numba's string formatting
utilities as they mature, and extending the diagnostic layer with
AST-structural pattern matching for improved coverage of nested unsupported
constructs.

% ---- References -------------------------------------------------------------
\begin{thebibliography}{15}

\bibitem{lam2015numba}
S.~K.~Lam, A.~Pitrou, and S.~Seibert,
``Numba: A LLVM-based Python JIT compiler,''
in \textit{Proc. 2nd Workshop LLVM Compiler Infrastructure in HPC (LLVM-HPC)},
2015, pp.~1--6.

\bibitem{liu2017performance}
Y.~Liu, S.~Huang, and C.~Ding,
``A performance study of LLVM-based just-in-time compilation for dynamically
typed languages,''
\textit{IEEE Trans. Parallel Distrib. Syst.},
vol.~28, no.~7, pp.~1873--1887, Jul. 2017.

\bibitem{oliphant2007python}
T.~E.~Oliphant,
``Python for scientific computing,''
\textit{Comput. Sci. Eng.},
vol.~9, no.~3, pp.~10--20, May 2007.

\bibitem{behnel2011cython}
S.~Behnel, R.~Bradshaw, C.~Citro, L.~Dalcin, D.~S.~Seljebotn, and K.~Smith,
``Cython: The best of both worlds,''
\textit{Comput. Sci. Eng.},
vol.~13, no.~2, pp.~31--39, Mar. 2011.

\bibitem{lehtosalo2022mypyc}
J.~Lehtosalo \textit{et al.},
``mypyc: Compiling Python to C using mypy type information,''
in \textit{Proc. Python in Science Conf. (SciPy)},
2022, pp.~44--50.

\bibitem{rigo2006pypy}
A.~Rigo and S.~Pedroni,
``PyPy's approach to virtual machine construction,''
in \textit{Proc. Dynamic Languages Symp. (DLS)},
2006, pp.~944--953.

\bibitem{cosenza2022numbadpex}
I.~Cosenza, N.~Gowanlock, B.~Juurlink, and P.~Millet,
``numba-dpex: A data-parallel extension for Numba,''
\textit{SoftwareX},
vol.~18, p.~101040, 2022.

\bibitem{zhang2024python}
T.~Zhang, X.~Liu, and Y.~Luo,
``Python meets JIT compilers: A simple implementation and a comparative
evaluation,''
\textit{Softw. Pract. Exp.},
vol.~54, no.~3, pp.~412--431, 2024.

\bibitem{psf2023ast}
Python Software Foundation,
``\texttt{ast} --- Abstract Syntax Trees,''
Python~3.11 Documentation, 2023.
[Online]. Available: \url{https://docs.python.org/3/library/ast.html}

\bibitem{vanrossum2014pep484}
G.~van~Rossum and P.~J.~Eby,
``PEP~484 --- Type Hints,''
Python Enhancement Proposals, 2014.
[Online]. Available: \url{https://peps.python.org/pep-0484/}

\bibitem{flatt2002macros}
M.~Flatt,
``Composable and compilable macros: You want it when?,''
in \textit{Proc. Int. Conf. Functional Programming (ICFP)},
2002, pp.~72--83.

\bibitem{nethercote2007valgrind}
N.~Nethercote and J.~Seward,
``Valgrind: A framework for heavyweight dynamic binary instrumentation,''
in \textit{Proc. ACM SIGPLAN Conf. PLDI},
2007, pp.~89--100.

\bibitem{bolz2009tracing}
C.~F.~Bolz-Tereick and A.~Rigo,
``Tracing the meta-level: PyPy's tracing JIT compiler,''
in \textit{Proc. ICOOOLPS},
2009, pp.~18--25.

\bibitem{rauchwerger1999lrpd}
L.~Rauchwerger and D.~Padua,
``The LRPD test: Speculative run-time parallelization of loops with
privatization and reduction parallelization,''
\textit{IEEE Trans. Parallel Distrib. Syst.},
vol.~10, no.~2, pp.~160--180, Feb.~1999.

\bibitem{seljebotn2009fast}
D.~S.~Seljebotn,
``Fast numerical computations with Cython,''
in \textit{Proc. Python in Science Conf. (SciPy)},
2009, pp.~15--22.

\end{thebibliography}

% ---- Author bio -------------------------------------------------------------
\begin{IEEEbiographynophoto}{Rodrigo Martinez Pinto}
is a software engineer focused on high-performance Python and cloud
cost-optimization tooling. He has contributed bug fixes and features to the
Numba upstream compiler (pull requests \#10880, \#10881, \#10882 accepted into
the numba/numba repository) and is the author of the open-source package
\texttt{cloudsealed-jit}. Contact: contact@cloudsealed.com
\end{IEEEbiographynophoto}

\end{document}
